% GENERATED FILE. Edit canonical lessons, then run npm run generate:books.
% canonical-source-hash: fnv1a64:819f7a04887dbe5e
% canonical-lessons: ES-C45-nos, ES-C45-nosotros

\chapter{Nos --- Us, and the We the Verb Already Says}
\label{ch:pronoun-nos}
\hlchaptermodality{88}

\begin{chapteropening}
\textbf{By the end of this chapter:} \emph{I can say nos quieren, and say what nosotros is made of.}

Complete ES-C45-nos: produce nos quieren and name the verb ending that already says we.
\end{chapteropening}

\section[nos]{nos — \textquotedblleft{}us\textquotedblright{}}

\label{lesson:ES-C45-nos}



\begin{quote}
Say \textquotedblleft{}you love me.\textquotedblright{} (\emph{Me quieres}.) And \textquotedblleft{}I love you.\textquotedblright{} (\emph{Te quiero}.) Both of those are about \textbf{one} person on the receiving end. Here is more than one.

\emph{Podemos} keeps the plain stem; \emph{pueden} breaks it again. Say both.
\end{quote}

\begin{grammarlens}[title={Grammar Lens: the third person in the room}]
\begin{quote}
\textbf{nos} — \emph{us}.
\end{quote}

\begin{quote}
\emph{\textbf{Nos} quieren.} — They love \textbf{us}.
\end{quote}

Same slot as always, in front of the verb. \emph{Me} is one person speaking; \emph{nos} is that person and somebody else.

\begin{quote}
\emph{Me quieren.} — They love me. \emph{Nos quieren.} — They love us.
\end{quote}

One letter of difference, and the difference is how many people you are speaking for.
\end{grammarlens}

\subsection*{Your turn}
Say these aloud:

\begin{itemize}
\raggedright
  \item \textquotedblleft{}Me quieren\textquotedblright{} then \textquotedblleft{}Nos quieren\textquotedblright{}
  \item \textquotedblleft{}Te quiero\textquotedblright{} then \textquotedblleft{}Nos quieren\textquotedblright{} — one person, then a group
\end{itemize}

\begin{tcolorbox}[breakable,colback=teal!4,colframe=teal!35!black,title={Before you move on}]
Say \textquotedblleft{}they love us.\textquotedblright{} (\emph{\textbf{Nos quieren}}.) And which slot did \emph{nos} take? (The one in \textbf{front of the verb}.) Next: where \emph{nos} came from, and the verb ending that has been saying \textquotedblleft{}we\textquotedblright{} without it.
\end{tcolorbox}
\section[nosotros]{nosotros — the \textquotedblleft{}we\textquotedblright{} the verb already says}

\label{lesson:ES-C45-nosotros}



\begin{quote}
Say \textquotedblleft{}they love us.\textquotedblright{} (\emph{Nos quieren}.)

\emph{Nos} is \textquotedblleft{}us.\textquotedblright{} Today it becomes \textquotedblleft{}we\textquotedblright{} — and you will find you have been saying \textquotedblleft{}we\textquotedblright{} for a long time without it.
\end{quote}

\begin{cousinweb}
\emph{Nos} is Latin \textbf{nōs}, which meant both \emph{we} and \emph{us} — the language had not yet split them into two words.

You have been saying \textquotedblleft{}we\textquotedblright{} without it. The \emph{-mos} ending on \emph{hablamos}, \emph{comemos}, \emph{tenemos} already means \textbf{we} — it is the verb's own \textquotedblleft{}we\textquotedblright{} ending, Latin \textbf{-mus} (\emph{amāmus}, \textquotedblleft{}we love\textquotedblright{}), a different word from \emph{nōs} that happens to mean the same person. That is why \emph{hablamos} alone is a whole sentence, and the pronoun is extra.

And the word for \textquotedblleft{}we\textquotedblright{} itself:

\begin{quote}
\textbf{nosotros} — \emph{we}, literally \textbf{nos} + \textbf{otros}, \textquotedblleft{}we \textbf{others}.\textquotedblright{}
\end{quote}

Spanish added \emph{otros} to the plain \emph{nos} somewhere in the Middle Ages, to make clear the speaker meant \emph{this} group and not the person being addressed. Catalan (\emph{nosaltres}) and Occitan (\emph{nosautres}) did the same; standard French and Italian did not, and in Spanish the seam is still visible in the word.
\end{cousinweb}

\subsection*{Your turn}
Say these aloud:

\begin{itemize}
\raggedright
  \item \textquotedblleft{}hablamos\textquotedblright{} — \textquotedblleft{}we speak,\textquotedblright{} with no pronoun at all
  \item \textquotedblleft{}nosotros hablamos\textquotedblright{} — \textquotedblleft{}we speak,\textquotedblright{} said twice
\end{itemize}

\begin{tcolorbox}[breakable,colback=teal!4,colframe=teal!35!black,title={Before you move on}]
What is \emph{nosotros} made of? (\emph{\textbf{Nos}} + \emph{\textbf{otros}} — \textquotedblleft{}we others.\textquotedblright{}) And which ending already says \textquotedblleft{}we\textquotedblright{} without it? (\emph{\textbf{-mos}}, from Latin \emph{-mus}.) Next: the plural of \emph{lo} and \emph{la}, and you will not need to be told what they are.
\end{tcolorbox}
